\section[几何类域论]{几何类域论（3 月 9 日）}
\srcpage{35}

今天讨论几何类域论，并开始（可能延续到下一讲）讨论它如何是一种范畴意义上的谱分解。

设 $\Sigma$ 是黎曼面。则 $\Pic(\Sigma)$ 落在短正合列
\[
0\longrightarrow \mathrm{Jac}(\Sigma)\longrightarrow \Pic(\Sigma)\longrightarrow \Z\longrightarrow 0
\]
之中，而 $\mathrm{Jac}(\Sigma)$ 是环面：它是某个向量空间对 $H_1(\Sigma;\Z)$ 的商。因此由 Poincaré 对偶，
$\pi_1(\mathrm{Jac}(\Sigma))\cong H_1(\Sigma;\Z)$。\footnote{稍细一点说，$\pi_1(\mathrm{Jac}(\Sigma))$
典范地就是 $H_1(\Sigma;\Z)$，Poincaré 对偶再把它与 $H^1(\Sigma;\Z)$ 联系起来。}
也就是说，Jacobian 的基本群与 $\Sigma$ 自身的基本群密切相关：
\[
\pi_1(\Sigma)^{\mathrm{ab}}\cong \pi_1(\mathrm{Jac}(\Sigma)).
\]
对阿贝尔群 $A$，记 $\mathrm{Loc}_A(\Sigma)=\{\pi_1(\Sigma)\to A\}/\cong$。映射 $\pi_1(\Sigma)\to A$
都经过交换化，而交换化就是 $\pi_1(\mathrm{Jac}(\Sigma))$。所以对这些局部系统、或 $\Sigma$ 的阿贝尔覆叠，
我们尽可以用 $\mathrm{Jac}(\Sigma)$：局部覆叠的同构类一一对应。我们特别关心 $A=\C^\times$ 的情形，
此时
\[
\mathrm{Loc}_A(\Sigma)=\Hom\bigl(\widehat{A},\mathrm{Loc}_1(\Sigma)\bigr),
\qquad \mathrm{Loc}_1(\Sigma)\cong \mathrm{Loc}_1(\mathrm{Jac}(\Sigma)).
\]

若取基点 $e\in\Sigma$，就有映射 $\Sigma\to\mathrm{Jac}(\Sigma)$，$x\mapsto x-e$，
它把 $\Pic^1$ 与 $\Pic^0=\mathrm{Jac}$ 等同起来。于是得到 Abel–Jacobi 映射
\[
\mathrm{AJ}^*:\mathrm{Loc}_1(\Pic^1)\xrightarrow{\ \sim\ }\mathrm{Loc}_1(\Sigma).
\]
这等于说：$\Pic^1$ 的 $A$-覆叠与 $\Sigma$ 的 $A$-覆叠是同一回事。

（TODO：这跟 Albanese 性质有关（TODO：我在这里漏讲了一点）。）若 $A$ 是任意阿贝尔群概形，
则映射 $\Sigma\to A$ 都经过 Abel–Jacobi 映射 $\Sigma\to\Pic^0(\Sigma)$。经典情形归功于
Lang–Rosenlicht；现代还有更一般的表述，例如 Justin Campbell 对叠给出的漂亮推广 [Cam17]。
也就是说，$\mathrm{AJ}^*$ 是同构
\[
\mathrm{Map}(\Sigma,A)\xleftarrow{\ \sim\ }\Hom_{\mathrm{Grp}}(\Pic(\Sigma),A).
\]
可以把它理解成一句话：「$\Pic(\Sigma)$ 是由 $\Sigma$ 自由生成的阿贝尔群概形。」

类似的事情在更一般的情形也成立：若 $X$ 是光滑射影簇，则有 Albanese 簇 $\mathrm{Alb}(X)$，
它也是阿贝尔群概形，可由 $X$ 上的 $0$-循环模去一个等价关系构造出来。
（即 $0$-循环是点的 $\Z$-线性组合；若它们落在同一个 $\PP^1$ 族里，就把它们等同。）
\srcpage{36}
$\mathrm{Alb}(X)$ 有万有性质：从 $X$ 到任意阿贝尔群概形的映射都经过 $\mathrm{Alb}(X)$。
在高维情形，$\Pic(X)$ 的定义方式不同（它关乎线丛）。

\begin{remark}
这与同伦论中一条重要定理有关，即 Dold–Thom 定理：由带基点空间 $X$ 作出的自由交换幺半群
$\Sym^\infty(X,*)$，其同伦群同构于 $X$ 的同调群。
\end{remark}

对我们这些黎曼面也可以做类似的事。例如，如前所见 $\Sigma\hookrightarrow\Pic^1$，
而 $\Sym^2\Sigma$ 映到 $\Pic(\Sigma)$，其中 $\Sym^k(\Sigma)$ 表示 $\Sigma$ 中 $k$ 个互异点构成的空间。
到 $\Pic^2$ 的映射把 $(x,y)$ 送到 $x+y$。在 $\Sym^d(\Sigma)$ 上照做，就落进 $\Pic^d(\Sigma)$：
这个映射保持次数。Abel–Jacobi 映射还可以写成
\[
\mathrm{AJ}:\Sym^*(\Sigma)\longrightarrow \Pic(\Sigma).
\]
$\Sym^*(\Sigma)$ 里那些点的有限组合，正是 $\Div(\Sigma)$ 中的有效除子；
后者可以等同于「全纯线丛 $L$ 连同截面 $\sigma$」再模去截面的数乘。Abel–Jacobi 映射只是忘掉截面。
因此线丛 $L\in\Pic(\Sigma)$ 上的纤维是 $\PP\bigl(H^0(\Sigma,L)\bigr)$。

这个映射不是满射，例如大多数线丛没有全纯截面。但当次数 $d$ 足够大时，由 Riemann–Roch 定理，
它们就有截面。具体地，需要 $d>2g(\Sigma)+2$，此时纤维是 $\PP^{d-g}$。

这很好：射影空间是单连通的，所以当次数大于 $2g(\Sigma)-2$ 时，Abel–Jacobi 映射在 $\pi_1$
上诱导同构。这又给出关于线丛的结论：可以把它看成从 $\Sigma$ 生成的自由交换幺半群到曲线生成的
自由阿贝尔群的群完备化映射，这正是 Dold–Thom 的精神。

我们想把 $\mathrm{Loc}(\Pic)$ 与 $\mathrm{Loc}_1(\Sigma)$ 的关系弄得更清楚。具体地，
我们希望有一个等价
\[
\mathrm{Loc}(\Pic)\cong \QCoh\bigl(\mathrm{Loc}_1(\Sigma)\bigr).
\]
一个麻烦是 $\Sigma$
只看见 $\Pic(\Sigma)$ 的一个连通分支，所以上面的故事总得改动一下；几何类域论要做的正是这件事。

回忆域 $\F_q$ 上曲线 $C$ 的无分歧类域论：我们把秩一局部系统与 $\Pic(C)$ 的特征标等同起来，
给局部系统 $E\in\mathrm{Loc}_1(C)$ 配上特征标 $\chi_E$。「特征标」意味着
\[
\chi_E(L\otimes M)\simeq \chi_E(L)\cdot \chi_E(M).
\]
又因为 $\Pic$ 由所有 $x\in C$ 对应的 $\Z x$ 生成，我们还知道
\[
\chi_E\bigl(L(x)\bigr)=\chi_E\bigl(L\otimes \cO(x)\bigr)
=\chi_E\bigl(\cO(x)\bigr)\chi_E(L).
\]
也就是说，$\chi_E(L)$ 是「按 $x$ 扭一下」这个操作的特征函数，其特征值是 $\chi_E(\cO(x))$。
而后者正是局部系统 $E$ 在 Frobenius 上的取值。并且 $\chi_E(\cO(x))=\mathrm{AJ}(x)$。

这件事的几何版本，是考虑 $\Pic$ 上的特征标局部系统，即群 $\Pic$ 上特征标的一种类似物。
这里的「特征标」不是表示论意义下的特征，而是群对象到 $\G_m$ 的乘法特征的层论类比
\yizhu{在特征层（character sheaf）理论中，「特征标」指 $\G_m$ 的局部系统沿群的乘法结构拉回所得的层；本译文统一作「特征标」。}{I（本库 SY02-21「特征标层」）}。
我们的场景是任意域 $k$ 上的曲线 $C$，但心里想的是黎曼面。$\Pic$ 上的特征标局部系统，
应当满足乘性条件（即上面那两条关系）。

\begin{definition}
$\Pic$ 上的特征标局部系统 $\chi\in\mathrm{Loc}(\Pic)$ 是一个局部系统，连同同构
$\mu^*\chi\xrightarrow{\ \sim\ }\chi\boxtimes\chi$ 的数据，其中 $\mu$ 是 $\Pic$ 上的乘法映射。
（TODO：当时还有没有别的要求？）
\end{definition}

逐点地看，把「$\mu$ 拉回后变成什么」讲清楚，就是说
\[
\chi|_{L\otimes M}\cong \chi|_L\otimes \chi|_M,
\]
但更自然的说法是在族上说。另一种看法是：$\chi$ 是映射 $\Pic\to B\G_m$。
利用乘性可以证明：阿贝尔群上任何具备该性质的局部系统必为秩一。
因为 $\chi(0+a)=\chi(a)=\chi(0)\otimes\chi(a)$，所以 $\chi(0)\cong k$。

可以考虑总空间 $A_\chi\to A$；$\chi$ 是特征标这一条件，等价于 $A_\chi$ 是阿贝尔群。
特别地，此时它是扩张
\[
0\longrightarrow k^\times\longrightarrow A_\chi\longrightarrow A\longrightarrow 0.
\]
这个群总是阿贝尔的；可以谈论概形中的阿贝尔群对象（从函子观点的角度看，就是到 $\mathrm{Ab}$
的可表函子），而这个扩张落在阿贝尔群的链复形范畴里（TODO：我想是这样）。
\srcpage{37}

特征标局部系统是范畴化的特征标，在任意阿贝尔群上都能定义。我们特别关心 $\Pic(C)$，
其中 $C$ 是任意域 $k$ 上的曲线。正如普通特征标那样，可以问 $\chi(L(x))$ 与 $\chi(L)$
的关系，类比于前面的乘性关系。我们得到
\[
\chi\bigl(L(x)\bigr)\cong \chi\bigl(\cO(x)\bigr)\otimes \chi(L),
\]
并且 $\cO(x)=\mathrm{AJ}(x)$。

令 $L_\chi:=\mathrm{AJ}^*\chi\in\mathrm{Loc}_1(C)$。于是乘性性质说的是
TODO（这里漏了，抱歉！）。这件事可以在族上说：考虑
\[
\begin{tikzcd}
\Sigma\times\Pic \arrow[r,"\mathrm{AJ}\times\id"] \arrow[dr,"\mu_\Sigma"'] & \Pic\times\Pic \arrow[d,"\mu"] \\
& \Pic
\end{tikzcd}
\]
断言是 $\mu_\Sigma^*(\chi)\simeq L\boxtimes\chi$，这正是 Hecke 特征性质。有了这一点，
我们就能真正陈述几何类域论：映射 $\chi\mapsto L_\chi=\mathrm{AJ}^*\chi$ 给出 $\Pic$ 上的
特征标局部系统与 $\Sigma$ 上秩一局部系统之间的等价。这是群胚的等价，也被看成叠的等价。
证明足够抽象，因而对任意域都成立；在环上也有类似物。

这还告诉我们：$\Sigma$ 的阿贝尔覆叠并不等同于 $\Pic$ 的任意阿贝尔覆叠，
而是等同于那些作为 $\Pic$ 的扩张出现的覆叠。所以我们确实需要知道 $\Pic$ 的群结构。
换句话说，$\Sigma$ 上任何秩一局部系统，典范地延拓到 $\Sym^*\Sigma$ 之后，都会下降落到 $\Pic$。

\subsection{Deligne 对几何类域论的证明}

Deligne 的证明先把 $L$ 延拓到 $\Sym^*(\Sigma)$。$\Sym^d(\Sigma)$ 的元素是点的有限和，
于是令
\[
L^{(d)}\Bigl(\sum_{i=1}^{d}x_i\Bigr):=\bigotimes_{i=1}^{d}L(x_i),
\]
这确实定义了一个秩一局部系统。更形式化地说，可以把这件事说成：给 $L$ 配上线丛
$L^{\boxtimes d}\to\Sigma^d$，再对对称群取商（我想是在挖掉对角线之后），
取 $\pi:\Sigma^d\to\Sym^d(\Sigma)$，得到 $(\pi^*L^{\boxtimes d})^{S_d}$。

下一步要用到：当 $d$ 足够大时，映射 $\Sym^d\to\Pic^d$ 是一连通的，因为纤维是射影空间，
而射影空间单连通。这意味着 $L^{(d)}$ 下降到 $\Pic^d$，即 $L^{(d)}\cong \mathrm{AJ}^*(\chi_L)$，
其中 $\chi_L$ 是 $\Pic^d$ 上的某个局部系统。这一步用到射影空间的单连通性，需要几何。

总之，对足够大的次数（$d>2g-2$）我们构造出了 $\chi_L$，而 $\chi_L$ 满足乘性
（尽管是受限的乘性：只对足够大的次数成立）。接下来的主张是：存在唯一的延拓到整个 $\Pic$、
仍满足乘性的 $\chi_L$。因为任何除子都能写成两个大次数有效除子之差 $D=D_1-D_2$，
其中 $D_1,D_2\in\im(\mathrm{AJ})$ 且次数足够大。乘性提示我们这样延拓 $\chi_L$：
\[
\chi_L(D):=\chi_L(D_1)\otimes \chi_L(D_2)^{-1}.
\]
Bhatt 有一份笔记（一篇 Oberwolfach 报告）讲几何类域论，是这个证明的好参考。
Bhatt 还指出，同样的证明对分歧几何类域论也有效：取 $S\subset\Sigma$，
考虑 Abel–Jacobi 映射 $\Sigma\setminus S\to\Pic(\Sigma,S)$，后者是在 $S$ 附近可平凡化的线丛构成的群。
这里用到了非阿贝尔情形下我们无法企及的一点：$\Pic$ 与曲线 $\Sigma$ 非常接近。
一般的图景会有更多在这里平凡掉的对偶性，也会更复杂。

下次我们换一个角度讲类域论，把它与谱分解的范畴化版本联系起来。我们想把 $\Pic$ 上的局部系统
与 $\mathrm{Loc}$ 上的层之间的关系，理解成某种 Fourier 变换。我们要对 $\Pic$ 上所有局部系统作谱分解
（不只是特征标局部系统——正如普通 Fourier 变换里，我们可以试着把任意局部系统写成特征标局部系统的组合）。
在 B 侧，看 $\QCoh(\mathrm{Loc}_1(\Sigma))$，我们的「基」是 $\mathrm{Loc}_1(\Sigma)$ 的点。

但我们的「谱」概念是范畴化的。课程开头我们讨论过：谱是「给几何对象（例如空间）配函数环」这个函子的
某种伴随或逆。把范畴数提高一格，我们把空间 $X$ 送到 $X$ 上的层的对称幺半范畴；因此
\srcpage{38}
谱函子应当朝相反方向，把一个对称幺半范畴 $\cC$ 尽量好地送到一个空间，也就是说存在伴随
\[
\mathrm{Map}\bigl(-,\Spec(\cC)\bigr)=\Hom_{\otimes\text{-}\mathrm{Cat}}\bigl(\cC,\mathrm{Sh}(-)\bigr).
\]
在这个场景里，这就叫 Tannakian 重构。点睛之笔将是
\[
\Spec\bigl(\mathrm{Loc}(\Pic)\bigr)=\mathrm{Loc}_1(\Sigma),
\]
这是几何类域论的一个特别强的版本。

若要在正特征中做这件事，就有「要哪种层」的问题——系数取在哪个域上的向量空间层？
对函数，我们用的是 $\Ql$ 值函数（非典范地等同于 $\C$ 值函数），而这里（TODO：我想）
会有类似的东西：若 $C$ 是有限域上的曲线，就取 $\Ql$ 系数的层。

在所有这些情形里，Tannakian 的视角最容易推广到非阿贝尔情形。
